\documentclass[aps,pra,reprint,amsmath,amssymb,superscriptaddress]{revtex4-2}
\usepackage{graphicx}
\usepackage{booktabs}
\usepackage{amsmath}
\usepackage[colorlinks=true, linkcolor=blue, citecolor=blue, urlcolor=blue]{hyperref}
\usepackage{xcolor}
\usepackage{orcidlink}

% ---------------------------------------------------------------------------
% Every number in this paper is \input from the repository's build, never typed
% by hand; see paper/README.md. Inline `% Claim N' comments map each result to
% its row in the repository's CLAIMS.md.
\input{../../build/latex/numbers.tex}

\graphicspath{{../../build/figures/}}

\begin{document}

\title{Gravity sourced by records: a non-signalling semiclassical law,\\
its predictions, and the limit on testing it}

\author{David T. Swanson~\orcidlink{0009-0003-0580-6476}}
\noaffiliation

\date{\today}

\begin{abstract}
I propose that gravity couples to a quantum system's state conditioned on the
physical records of its position that actually exist, whether or not they are
ever read, rather than to its full wavefunction or to what an observer has
measured. The law is stated in density-matrix language, provably does not
signal, and has one closed form selected by requiring that independent records
compose multiplicatively. It predicts no gravitational entanglement between two
unrecorded masses, with no noise added to enforce that; and no self-gravity
frequency shift for any mechanical oscillator above its motional ground state.
I verify these properties by exact simulation, with controls, and then ask
what a real experiment could see. On every platform considered, the records
that suppress the predicted signal also broaden the spectral feature that would
carry it, independent of the oscillator's mass; this limit is specific to the
present law, and a measurement-conditioned alternative evades it. A weaker
bound applies to the whole class: any spectroscopic search for a conditioned
second resonance needs an observation time of order $\pi Q/\omega_0$, a
special case of a bound derived for measurement-conditioned
Schr\"odinger--Newton gravity, and consistent with the run time reported by a
recent torsion-balance experiment.
\end{abstract}

\maketitle

% =============================================================================
\section{Introduction}

Whether gravity couples to quantum matter as a quantized field or through some
semiclassical channel is open, and if the latter, the open question sharpens:
semiclassical gravity needs a source, and \emph{which state} of a quantum
system that source is drawn from is itself a choice with physical
consequences. Two answers are on the table. Unconditional
Schr\"odinger--Newton gravity sources the field from a system's full
wavefunction, whatever branch it occupies and whatever records of it exist;
the resulting nonlinear, nonlocal theory is not excluded by any experiment
performed to date, but is argued on general grounds to permit
superluminal signalling~\cite{EppleyHannah1977} and is in direct tension with
the one experiment built to distinguish it from the ordinary, linear quantum
expectation value~\cite{PageGeilker1981}. Measurement-conditioned proposals
instead source the field from what an experimenter's own readout has
resolved~\cite{Yang2013,Helou2017,LiuEtAl2023}; this avoids sourcing
by an unmeasured superposition, but makes the source depend on a choice --- what to measure, and when.

This paper takes a third option: condition the source not on the full
wavefunction, and not on what has been measured, but on whatever physical
\emph{records} of a system's position exist, independent of whether anyone
ever reads them. A record is a copy --- a correlation with the environment, a
latch in an apparatus, a decohered pointer state --- and its existence is a
fact about the world, not about an observer's choice. I show that the
resulting law cannot signal (Sec.~\ref{sec:law}), and derive its
consequences for two
canonical scenarios: two masses held in spatial superposition with no record
of either taken (Sec.~\ref{sec:predictions}), and a single oscillator's
resonance spectrum, where the \emph{same} source rule makes a sharp,
nearly parameter-free prediction depending only on whether the oscillator's
own thermal state has been reduced to its motional ground state.

Section~\ref{sec:checks} verifies the law's properties by exact simulation in
a small model, with controls showing that the simulation distinguishes the law
from each alternative. Section~\ref{sec:realistic} extends the law from two
branches to a continuously decohered mechanical oscillator and evaluates it on
three published platforms. The result is negative: the mechanism that makes
the predicted signal small also makes the spectral feature carrying it broad,
and the two do not separate. Section~\ref{sec:why} shows why this limit is
specific to the present law, and Sec.~\ref{sec:mintime} identifies the weaker
limit common to every proposal of this kind. Section~\ref{sec:related}
compares the law with other proposals, and Sec.~\ref{sec:limits} lists the
limitations of the analysis.

I present the law as a hypothesis with computable consequences, and identify
what would test it.

% =============================================================================
\section{The law}
\label{sec:law}

\subsection{A record, in density-matrix language}

Let a quantum system occupy two distinguishable branches, labelled $A$ and
$B$ --- the two arms of an interferometer, the two positions of a levitated
mass, the two outcomes a decay could leave behind. A \emph{record} of which
branch is a correlation with some other degree of freedom, a register, whose
two conditional states $\lvert r_A\rangle$ and $\lvert r_B\rangle$ overlap at
\begin{equation}
  \eta \equiv \lvert \langle r_A | r_B \rangle \rvert \in [0,1].
\end{equation}
$\eta=1$ is no record at all (the register carries no information about which
branch occurred); $\eta=0$ is a perfect record. This is not a new
quantity: $\eta$ is exactly the system's own decoherence factor, the
off-diagonal element of its reduced density matrix in the branch basis,
normalized to the branch populations. For an exactly balanced
two-branch state,
\begin{equation}
  \eta = \frac{\lvert \rho_{AB}\rvert}{\sqrt{\rho_{AA}\rho_{BB}}}
       = \sqrt{2\,\mathrm{Tr}\,\rho^2 - 1},
  \label{eq:eta-purity}
\end{equation}
the second equality following from $\mathrm{Tr}\,\rho^2 = (1+\eta^2)/2$ for a
two-branch state at equal weight. Equation~\eqref{eq:eta-purity}'s
right-hand side is defined for \emph{any} state, not only a balanced
two-branch one, and this is what lets Sec.~\ref{sec:realistic} extend the law
to a real, continuously decohering oscillator with no branches in it at all.

The quantity $\eta$, and hence everything that follows, is a property of the
system's own \emph{reduced} state $\rho$, read in whatever basis locality
fixes. It says nothing about whether the register has been read, by whom, or
in what basis a distant party chooses to read it; Appendix~\ref{app:nosignal}
uses this to prove that the law does not signal.

\subsection{The source}

The gravitational source is the record-conditioned density $n_r$, built to
interpolate between two fixed points: with no record ($\eta=1$), $n_r$ is the
ordinary quantum-mechanical mean density, $n_r = (n_A+n_B)/2$; with a perfect
record ($\eta=0$), $n_r$ is the localized branch alone, $n_r = n_A$ (for the
branch-$A$ description). At intermediate $\eta$, two forms interpolate:
\begin{align}
  \text{(i)}\quad   n_r^A &= \eta\,\tfrac{n_A+n_B}{2} + (1-\eta)\,n_A, \\
  \text{(ii)}\quad  n_r^A &= (1-P)\,n_A + P\,n_B,\qquad
    P = \tfrac{1-\sqrt{1-\eta^2}}{2}.
\end{align}
Writing $f(\eta)$ for the weight the law places on the branch a given
description does \emph{not} occupy, normalized to its no-record value of
one half, the two forms are
\begin{equation}
  f_{\mathrm{(i)}}(\eta) = \eta,
  \qquad
  f_{\mathrm{(ii)}}(\eta) = 1-\sqrt{1-\eta^2}.
  \label{eq:f-forms}
\end{equation}
Both satisfy $f(1)=1$ and $f(0)=0$, are continuous and monotonic on $[0,1]$,
and contain no free parameters. Figure~\ref{fig:f-eta} plots both, together
with the drift between two branches they produce in an exact numerical model;
the drift follows $f(\eta)$ to within $4$--$6\%$, the difference coming from
each branch's own wavepacket spreading. % Claim 1

\begin{figure}[t]
  \centering
  \includegraphics[width=\columnwidth]{check-1-f-eta.pdf}
  \caption{$f(\eta)$ for both closed forms, Eq.~\eqref{eq:f-forms}
    (left), and an exact numerical model's reading of it (right): the
    measured drift between two branches, relative to its $\eta=1$ value,
    tracks $f(\eta)$ to within a few per cent. At $\eta=0.5$, $f_{\mathrm
    {(i)}}=\fEtaHalfRuleI$ against a measured ratio of
    \etaDriftRatioHalfRuleI; $f_{\mathrm{(ii)}}=\fEtaHalfRuleII$ against
    \etaDriftRatioHalfRuleII. The two forms differ by a factor of
    \fEtaHalfRatio\ at $\eta=0.5$.}
  \label{fig:f-eta}
\end{figure}

\subsection{Properties, and the form adopted}

\emph{No signalling.} Both forms are functions of $\eta$ alone, and $\eta$ is
a function of the system's own reduced state, Eq.~\eqref{eq:eta-purity}.
No local operation on the register --- reading it, not reading it, applying
any unitary to it whatsoever --- can change the system's reduced state, so
neither form of the law can be used to signal
(Appendix~\ref{app:nosignal}). Numerically, the prediction is unchanged to
$\noSignallingWorstTwoDim$ under twenty random unitaries on a two-dimensional
register and to $\noSignallingWorstFourDim$ on a four-dimensional one, for
both forms. A law sourced instead by the outcome actually read changes by
$\outcomeSourcedSignal$ of its no-record value under the same test.
% Claims 5, 6

\emph{Composition.} Two independent records at overlaps $\eta_1,\eta_2$ are
read by the law as a single record at overlap $\eta_1\eta_2$, exactly, for
\emph{either} form (to $\compositionWorst$), because the law reads only the
system's reduced state, on which two independent records act multiplicatively
in their overlaps. % Claim 3

\emph{Multiplicativity.} A stronger condition is that $f$ itself be
multiplicative, $f(\eta_1)f(\eta_2)=f(\eta_1\eta_2)$. Form (i), $f(\eta)=\eta$,
satisfies it exactly (error $\ruleIMultiplicativeError$); form (ii) misses by
up to $\ruleIIMultiplicativeMiss$ in absolute terms. This is the only
criterion among these checks that separates the two forms; no signalling holds
for both. % Claims 4, 6
I take multiplicativity as the selection principle and adopt form (i),
$f(\eta)=\eta$, as the law in what follows. The principle is a requirement
imposed here, not a consequence of the law's construction (Sec.~\ref{sec:limits}).
Form (ii) is retained in the figures and tables as the alternative.

\emph{No added noise.} Unlike a genuinely classical gravitational channel,
which must add decoherence to the systems it couples as the price of not
signalling~\cite{Kafri2014}, this law's mean-field dynamics is a
\emph{deterministic} function of the state, with no stochastic term. The law
avoids signalling by conditioning on a basis-independent fact, how much record
exists, rather than by adding noise to wash out a dependence on what is read.
% Claim 9

% =============================================================================
\section{Predictions}
\label{sec:predictions}

\paragraph{P1: gravity without entanglement, and without added noise.}
Two masses, each held in a spatial superposition with no record of either
taken, interact only through their respective mean fields under this law.
Simulated exactly as two qubits, the state's negativity stays below
$\negativityMax$, zero to numerical precision, while a quantized two-body
coupling applied to the same initial state reaches a concurrence of
$\quantizedConcurrenceMax$ (Fig.~\ref{fig:twomass}). % Claim 8
The masses still attract: each accumulates a relative phase from the other's
mean field. What is absent is entanglement and, unlike a classical-channel
model~\cite{Kafri2014}, which also predicts no entanglement but adds noise,
any added noise. This is the
discriminator the gravitationally-induced-entanglement
proposals~\cite{BoseEtAl2017,MarlettoVedral2017} are built to test: if gravity
mediates entanglement between two otherwise-unconnected quantum systems, its
mediator must itself be quantum. Under this law, with no record taken, it
would not. The simulation is mean-field rather than a full master equation
(Sec.~\ref{sec:limits}).

\begin{figure}[t]
  \centering
  \includegraphics[width=\columnwidth]{check-4-two-mass.pdf}
  \caption{Two unrecorded masses: negativity under this law's mean-field
    form (flat, below $\negativityMax$) against a quantized two-body coupling
    applied to the same initial state (reaching concurrence
    $\quantizedConcurrenceMax$), which serves as the control.}
  \label{fig:twomass}
\end{figure}

\paragraph{P2: no self-gravity shift above the ground state.}
For a single mechanical oscillator, the law's weight on the unconditioned
state, $f(\eta)$, is nonzero only where the oscillator's own thermal
occupation satisfies $\bar n < 1/2$ --- i.e., only in or near the motional
ground state --- because $\eta_{\mathrm{eff}} = \sqrt{2\,\mathrm{Tr}\,\rho^2-1}$
with $\mathrm{Tr}\,\rho^2 = 1/(2\bar n+1)$, and both factors are forced by
Eq.~\eqref{eq:eta-purity} (Sec.~\ref{sec:realistic}). For any oscillator with
$\bar n>1/2$, which includes all three platforms considered here, the
resonance sits exactly at
the bare mechanical frequency $\omega_0$, with \emph{no} self-gravity shift
at all. Unconditional Schr\"odinger--Newton, by contrast, shifts the
resonance to $\omega_q=\sqrt{\omega_0^2+\omega_{\mathrm{SN}}^2}$ regardless
of the oscillator's thermal state. Unconditional Schr\"odinger--Newton
gravity is not excluded by any published experiment (Sec.~\ref{sec:related}),
so this is a live discriminator.

\paragraph{P3: record dependence in the ground state.}
Only for an oscillator cooled below its motional ground-state occupation
threshold does the record weight $f(\eta)$ depart from zero, and with it the
signatures that distinguish an unread record from a read one, and a held
record from an erased one. Section~\ref{sec:realistic} shows these signatures
are out of reach on every platform considered: the one platform pure enough to
have $f\neq0$ is also the one where the predicted signal is smallest relative
to the noise.

\begin{table*}[t]
  \centering
  \caption{Discriminators: this law (both forms) against quantized gravity,
    standard quantum mechanics, and the unconditional and
    measurement-conditioned Schr\"odinger--Newton alternatives, across four
    scenarios. Dashes mark combinations not modelled here; see the
    footnotes.}
  \label{tab:discriminators}
  \footnotesize
  \input{../../build/latex/tables/discriminators-full.tex}
  \vspace{2pt}
  {\footnotesize
  $^{\dagger}$Exact two-qubit simulation, mean-field form (Fig.~\ref{fig:twomass}).\quad
  $^{\ddagger}$Zero by construction: no coupling term exists between the masses.\quad
  $^{a}$Not modelled here: neither unconditional nor measurement-conditioned
  Schr\"odinger--Newton has a two-mass mechanism implemented here; the
  classical-channel answer is no entanglement with added
  noise~\cite{Kafri2014}.\quad
  $^{b}$No two-body coupling partner exists in a single-oscillator spectrum;
  trivially identical to standard quantum mechanics in this scenario.\quad
  $^{c}$No register is implemented for these laws here.\quad
  $^{d}$Unread means nothing has yet been read, so the measurement-conditioned
  source has not conditioned on anything --- identical to the no-record case
  here.
  }
\end{table*}

Table~\ref{tab:discriminators} collects the four scenarios across the laws
compared.

% =============================================================================
\section{Exact checks}
\label{sec:checks}

The law and its three comparison laws (unconditional Schr\"odinger--Newton,
a measurement-conditioned alternative that sources by the outcome actually
read, and standard quantum mechanics with no self-gravity term) are
implemented in one shared code path, each differing only in what the
conditioned density $n_r$ is set to. Two models carry the checks: an exact
split-step propagation of two branch wavefunctions, each in its own
record-conditioned potential, on a few hundred grid points; and exact linear
algebra on a register Hilbert space of at most a few hundred dimensions.
Nothing is sampled or fitted.

\begin{table}[t]
  \centering
  \caption{The nine tests and their results.}
  \label{tab:tests}
  \footnotesize
  \input{../../build/latex/tables/tests-summary.tex}
\end{table}

Table~\ref{tab:tests} summarizes the nine tests. Three controls show that the
simulation distinguishes the laws: unconditional Schr\"odinger--Newton is flat
in $\eta$, as it must be; the quantized two-body coupling generates
entanglement where the law's mean-field form does not
(Fig.~\ref{fig:twomass}); and the outcome-sourced alternative differs from
the law when a record goes unread. % Claims 8, 15, 16

\begin{figure}[t]
  \centering
  \includegraphics[width=\columnwidth]{check-2-no-signalling.pdf}
  \caption{A distant party's choice of measurement basis on a perfect record:
    the law's prediction (both forms) is identical in either basis; an
    outcome-sourced alternative's differs by $\outcomeSourcedSignal$ of its
    no-record value between them, and so signals.}
  \label{fig:nosignal}
\end{figure}

Figure~\ref{fig:nosignal} shows the no-signalling check directly.

For the conservation test (T9), the conserved quantity is the functional that
generates the dynamics,
\begin{equation}
  E = T_A+T_B+\tfrac{\alpha}{2}\big(\langle n_AKn_A\rangle+\langle
  n_BKn_B\rangle\big)+\beta\langle n_AKn_B\rangle,
\end{equation}
with $\alpha,\beta$ the law's two coefficients. (Summing each branch's
potential against the other branch's density instead double-counts the
self-interaction.) The worst drift is $\energyDriftWorst$, and it falls by a
factor of $\energyDriftQuarterRatio$ when the time step is quartered, as
expected for first-order integrator error. Across a change of record the
energy changes by $\energyDriftAcrossErasure$, because changing the record
changes the Hamiltonian; testing conservation across such a change requires
the full mass-plus-register system, which this reduced model does not include
(Sec.~\ref{sec:limits}).

% =============================================================================
\section{Under realistic conditions}
\label{sec:realistic}

\subsection{The adapter}

The checks above use idealized two-branch states. A real test uses a massive
mechanical oscillator, continuously decohered by its thermal environment and
continuously read out by light. Extending the law to this setting takes three
steps.

\emph{How much record exists.} Equation~\eqref{eq:eta-purity}'s
right-hand side, $\eta_{\mathrm{eff}}=\sqrt{2\,\mathrm{Tr}\,\rho^2-1}$, is
read off the oscillator's \emph{unconditioned} steady-state covariance --- the
state whose mixedness the environment's records caused.

\emph{What the source is.} The law's own interpolation is applied with the
unconditioned density playing the no-record role and the state
\emph{conditioned on every record that exists} (the continuous-measurement
filtering solution, at unit efficiency on every decoherence channel) playing
the perfect-record role. This reduces exactly to the two-branch law in the
two-branch limit. % Claims 11, 12

\emph{What it does.} In the regime where the source's spread is small against
the material's lattice localization length, the self-gravity term is
harmonic, and the conditioned source produces \emph{two} resonances: the
conditional mean oscillates at $\omega_c=\sqrt{\omega_0^2+f\,\omega_{\mathrm
{SN}}^2}$, and the fluctuation about it, which carries the conditional
variance of whatever the law conditions on, sits at
$\omega_q=\sqrt{\omega_0^2+\omega_{\mathrm{SN}}^2}$.

This construction extends the law beyond the two-branch states for which it
is stated. Appendix~\ref{app:adapter} gives the filtering equations and an
independent check showing that the conclusions below do not depend on the
extension. % Claim 20

\subsection{Three platforms}

\begin{table}[t]
  \centering
  \caption{The three platforms, drawn from the published literature, with
    assumed values marked.}
  \label{tab:platforms}
  \input{../../build/latex/tables/platforms.tex}
\end{table}

Table~\ref{tab:platforms} lists the three platforms: a
superconducting osmium microdisc in a Paul trap~\cite{G16}, the published
proposal's own choice of material and trap frequency; a milligram torsion
balance~\cite{W21}, the lowest mechanical frequency measured to date; and a
ground-state-cooled levitated nanosphere~\cite{D20}, the only platform of the
three whose thermal occupation is below one half, and so the only one where
$f(\eta)$ is nonzero at all. This sphere is silica, for which no
crystal-regime self-gravity constant is published; silicon's is used instead, which
overstates the effect and is therefore conservative for the negative
conclusion below (Appendix~\ref{app:platforms}). Parameters not taken from
their sources are marked as assumed and swept over many orders of magnitude.

\subsection{The result: a shape failure}

\begin{figure}[t]
  \centering
  \includegraphics[width=\columnwidth]{realistic-3-spectra.pdf}
  \caption{What the detector sees on the osmium platform: this law's
    prediction, the measurement-conditioned alternative's, and standard
    quantum mechanics' main resonances coincide; only unconditional
    Schr\"odinger--Newton's is shifted (left). The excess over standard
    quantum mechanics at the offset where a resolved second peak would sit,
    relative to each law's own linewidth there (right): the excess is
    present but spread far wider than the offset itself.}
  \label{fig:spectra}
\end{figure}

On every platform, the law's conditioned second resonance is broader than its
offset from the main line by four to eight orders of magnitude
(Table~\ref{tab:platforms}), % Claim 14
because the environment's own
records simultaneously suppress the predicted signal's \emph{weight} and set
the resonance's \emph{width}, through the same localization rate --- visible directly in Fig.~\ref{fig:spectra}. The obstacle is the line shape,
not sensitivity: no integration time resolves a feature far broader than its
own offset.

\begin{figure}[t]
  \centering
  \includegraphics[width=\columnwidth]{realistic-5-mass-independence.pdf}
  \caption{Resolvability as a function of oscillator mass, at fixed
    $\omega_0$, $Q$ and $T$: flat to $\massIndependenceSpread$ (floating-point
    noise) over fourteen orders of magnitude in mass. A heavier oscillator
    does not help.}
  \label{fig:massindep}
\end{figure}

The closed form underlying this, valid in the weak-measurement, small-shift
regime (Appendix~\ref{app:resolvability}),
\begin{equation}
  R \equiv \frac{\text{offset}}{\text{width}}
    = \frac{Q\,\hbar\,\omega_{\mathrm{SN}}^2}{8\,\omega_0\,k_BT},
  \label{eq:resolvability}
\end{equation}
does not depend on the oscillator's mass, which the full filtering model
confirms to $\massIndependenceSpread$ over fourteen orders of magnitude
(Fig.~\ref{fig:massindep}). % Claim 17
A larger mass therefore does not help. $R$ improves
with $Q/T$: osmium needs $Q/T\sim\qtOsmiumMilliHz\,\mathrm{K^{-1}}$ at a
millihertz to reach $R=1$, or $\qtOsmiumTenHz\,\mathrm{K^{-1}}$ at ten hertz.

\begin{figure}[t]
  \centering
  \includegraphics[width=\columnwidth]{realistic-6-q-window.pdf}
  \caption{The feasible region in $Q$ and $\omega_0$, at four temperatures: a
    \emph{window}, not a threshold. Below the lower curve the second
    resonance is an unresolved shoulder; above the upper curve, it is
    resolved but too narrow for a periodogram to gather enough independent
    samples within a month. The window closes from both sides as $\omega_0$
    rises.}
  \label{fig:window}
\end{figure}

A region meeting all three feasibility conditions (resolvability of one, and
separation from each alternative within a month) exists: osmium at
$1\,\mathrm{mHz}$, $Q=10^{12}$ and $1\,\mathrm{mK}$ gives $R=\cornerR$. It
is a window rather than a threshold (Fig.~\ref{fig:window}): % Claim 18
far enough above the floor, the resonance becomes too narrow for a
periodogram to gather enough independent samples in a month, so raising $Q$
eventually hurts.

\subsection{The observable}

Modulating an ancilla's coupling and watching the raw spectrum does not
isolate the effect: the power spectral density at the probe frequency moves by
up to $\qmLockinRaw\times$ its baseline under standard quantum mechanics,
which has no self-gravity term, because the ancilla's backaction heats the
oscillator under every law. The second resonance's \emph{weight}, separated
from this heating, passes every control. A protocol of this kind must use
that heating-nulled observable. % Claims 15, 16

% =============================================================================
\section{Why the limit is this law's}
\label{sec:why}

The measurement-conditioned alternative does not face the same wall. Its
second resonance's half-width is set by $\eta_{\det}\Lambda_{\mathrm{ro}}$,
the readout's own monitoring rate, which the experimenter sets and so can
tune down to the imprecision floor, $\Gamma_2=2\gamma$. This law's resonance
width is set instead by $\Lambda_{\mathrm{tot}}$, the environment's own total
record rate, which nobody chooses, giving $\Gamma_2=2\gamma(2\bar n+1)$ --- larger by the thermal occupation factor, $\thermalOccupationFactorOsmium$ on
the osmium platform and $\thermalOccupationFactorTorsion$ on the torsion
balance. % Claim 22
Scanning the readout strength over fifty-six orders of magnitude in the full
model, the measurement-conditioned law clears the imprecision floor on the
osmium platform with resolvability up to $4.8\times10^4$; the present law
clears it on no platform.

The law is hard to test not because its signal is small, but because it is
sourced by records that no one can choose not to make.

% =============================================================================
\section{A minimum observation time for the whole programme}
\label{sec:mintime}

A spectral feature of half-width $\Gamma$ cannot be resolved in less than
about its own coherence time, $2\pi/\Gamma$. Every law with a conditioned
second resonance --- this one and the measurement-conditioned alternative
alike --- has $\Gamma_2\ge2\gamma$, forced by the imprecision floor every
readout carries (Sec.~\ref{sec:why}). Combining the two gives a bound
independent of which law, platform, or confidence level is in question:
\begin{equation}
  t_{\min} = \frac{\pi\,Q}{\omega_0},
  \label{eq:tmin}
\end{equation}
growing \emph{with} $Q$ --- this is the reason the window in
Fig.~\ref{fig:window} has a ceiling. For a published torsion-balance
experiment~\cite{Y25} ($Q\approx5\times10^4$ at $0.6\,\mathrm{mHz}$),
Eq.~\eqref{eq:tmin} gives $\tMinPublishedTorsionDays$ days from its quality
factor and resonance frequency alone; the experiment reports needing more than
$300$ days to reach steady state, in agreement in order of magnitude.

The relation underlying Eq.~\eqref{eq:tmin}, that resolving a spectral feature
takes a time set by its inverse width, is standard. Helou
\emph{et al.}~\cite{Helou2017} derive a sharper version for
measurement-conditioned Schr\"odinger--Newton gravity,
$\tau_{\min}(h)\approx27\,h^{-0.73}\times2/\gamma$ for a Lorentzian peak of
normalized height $h$ at ten per cent confidence (their Eq.~120).
Equation~\eqref{eq:tmin} is the weaker bound obtained by dropping the
dependence on confidence level and signal strength; in exchange, it applies to
any law with a conditioned second resonance. I note its generality and its
agreement with the run time reported in Ref.~\cite{Y25}.

% =============================================================================
\section{Relation to other proposals}
\label{sec:related}

\paragraph{Unconditional Schr\"odinger--Newton.} This law reduces to it
exactly with no record ($f(1)=1$), and departs from it wherever any record
exists. It remains unexcluded: the most sensitive published
test~\cite{Y25} reports no evidence for a semiclassical-gravity effect and
explicitly does not claim to exclude the equation.

\paragraph{Measurement-conditioned Schr\"odinger--Newton.} Sourced by what an
experimenter's readout has resolved rather than by records that
exist~\cite{Yang2013,Helou2017,LiuEtAl2023}. It differs from the present law
for unread records, and it avoids the resolvability limit of
Sec.~\ref{sec:why}.

\paragraph{Classical-channel and spontaneous-localization models.}
Kafri, Taylor and Milburn~\cite{Kafri2014} treat gravity between two
resonators as a genuinely stochastic classical measurement channel: like this
law, it cannot entangle two masses through gravity alone, but unlike this
law, it must add decoherence as the price of not signalling, since a
classical channel cannot transmit quantum correlations noiselessly. Tilloy
and Di\'osi~\cite{TilloyDiosi2016} instead source gravity from the state a
spontaneous-localization model has already localized, eliminating the
single-particle self-interaction and predicting an added decoherence term
tied to the localization model's own parameters. Both share the present law's premise, that gravity is sourced by what is
localized rather than by the full superposition, but replace conditioning on
existing records with an intrinsically stochastic mechanism.

\paragraph{Classical--quantum gravity.} Oppenheim~\cite{Oppenheim2023}
couples a genuinely classical spacetime to quantum matter through a
trace-preserving, completely positive master equation, forcing quantum
mechanics itself to become fundamentally stochastic. This modifies the
matter dynamics directly, a different scope than the present law, which
modifies only the source term.

\paragraph{Collapse models.} Continuous spontaneous localization at Adler's
parameters~\cite{BassiRMP2013} already heats a representative platform's
oscillator eight to eleven orders of magnitude faster than its own thermal
diffusion rate, excluding it on thermal grounds alone, before any
self-gravity question is asked. The Di\'osi--Penrose criterion
\cite{Diosi1989,Penrose1996} and Ghirardi--Rimini--Weber~\cite{GRW1986} are,
like continuous spontaneous localization, collapse mechanisms rather than
sourcing laws: they remove a superposition rather than conditioning a source
on its records.

\paragraph{Gravitationally-induced entanglement.} Bose \emph{et
al.}~\cite{BoseEtAl2017} and Marletto and Vedral~\cite{MarlettoVedral2017}
propose witnessing the gravitational field's quantum nature by whether two
masses in spatial superposition become entangled through their mutual
attraction. This is precisely the test P1 (Sec.~\ref{sec:predictions})
speaks to: under this law, with no record taken, they would not.

\paragraph{The case for quantizing gravity.} Eppley and
Hannah~\cite{EppleyHannah1977} argue on general grounds that a classical
field coupled to quantum matter must violate momentum conservation, the
uncertainty principle, or no-signalling; Page and Geilker~\cite{PageGeilker1981}
give the experimental half of the case, a Cavendish-type experiment whose
source mass's position is entangled with a radioactive decay, which always
deflects toward one branch or the other rather than the expectation value
unconditional semiclassical gravity predicts. Both bear on sourcing gravity
from the unconditioned mean, not on conditioning it on records.

% =============================================================================
\section{Limits}
\label{sec:limits}

\begin{enumerate}
  \item The selection between the two closed forms of $f$ rests on one
    criterion --- multiplicative composition --- not on a deeper principle. No
    signalling, in particular, does not distinguish them.
  \item The crystal-regime self-gravity frequency $\omega_{\mathrm{SN}}$, and
    the assumption that it enters the oscillator's frequency squared linearly
    in the law's own weight, are imported from the optomechanics literature,
    not derived from this law.
  \item The adapter's reduction to a thermal oscillator's state (rank above
    two) is an extrapolation past where the law is stated for two branches.
    An independent reading, at the material's lattice scale rather than
    through the state's purity, gives the same conclusion on both platforms
    where it matters, but the extension itself is not part of the law.
  \item The two-body entanglement result (P1) is mean-field, not a full
    continuous master equation with decoherence and backaction acting
    throughout.
  \item Every decoherence channel in the adapter is treated as resolving
    position exclusively; a bath that instead recorded momentum would need a
    different adapter, not built here.
  \item Classical force noise --- Casimir forces, patch potentials --- is
    not included in the realistic-conditions model. It would add to the
    total diffusion rate and so strengthen the negative feasibility
    conclusion, but it has not been shown to be negligible.
  \item Two of the three platforms (the torsion balance and the levitated
    sphere) carry an assumed mechanical quality factor, swept broadly in
    Sec.~\ref{sec:realistic} but quoted as single numbers in
    Table~\ref{tab:platforms}.
  \item The minimum-observation-time bound (Sec.~\ref{sec:mintime}) is a
    special case of a known result~\cite{Helou2017} and is compared with a
    single experiment.
\end{enumerate}

% =============================================================================
\section{Conclusion}

Gravity sourced by the records that physically exist of a quantum system's
position --- rather than by its full wavefunction or by what an observer has
measured --- is a law that does not signal, by construction, and that makes
two sharp predictions: no gravitational entanglement between unrecorded
masses, with no added noise; and no self-gravity frequency shift for any
oscillator above its motional ground state. Exact simulation with controls
confirms both. The signature that distinguishes the law from its
measurement-conditioned alternative is out of reach on every platform
considered: the records that suppress the signal also broaden the feature
that would carry it, independent of the oscillator's mass. Prediction P2 can
be tested with existing apparatus, since none of the platforms in
Table~\ref{tab:platforms} is near its motional ground state, and it separates
the law from unconditional Schr\"odinger--Newton gravity. Finally, the bound of
Sec.~\ref{sec:mintime} applies to every proposal of this kind: sharpening the
resonance to see a conditioned self-gravity signal lengthens the required
observation time in proportion.

\paragraph{Data and code availability.} Every number, table and figure in
this paper is generated by the \texttt{record-conditioned-gravity}
repository~\cite{rcgrepo}, where \texttt{make all} reproduces them from a
clean checkout in under an hour. The repository's \texttt{CLAIMS.md} maps each
result to the script and output that produce it.

\paragraph{Origin.} The law arose in work on a separate, unpublished theory
of relational quantum gravity. It is stated here independently, and nothing in
this paper relies on that theory.

\appendix

% =============================================================================
\section{The closed forms}
\label{app:closedforms}

Equation~\eqref{eq:f-forms}'s two forms follow directly from the two
interpolation rules for $n_r$ given in Sec.~\ref{sec:law}, by evaluating the
weight each places on the non-occupied branch and normalizing to the
no-record value; both are closed forms in $\eta$ alone, implemented as
\texttt{rcg.law.f\_rule\_i} and \texttt{rcg.law.f\_rule\_ii}.

% =============================================================================
\section{No signalling}
\label{app:nosignal}

\textbf{Claim.} For any $f$ depending only on the reduced state $\rho$ of the
system under study, no local operation on a register entangled with that
system can change the law's prediction.

\textbf{Proof sketch.} The reduced state is $\rho=\mathrm{Tr}_R\,\lvert
\Psi\rangle\langle\Psi\rvert$, tracing out the register $R$. For any unitary
$U_R$ acting on $R$ alone, $(\mathbb{1}\otimes U_R)\lvert\Psi\rangle\langle
\Psi\rvert(\mathbb{1}\otimes U_R)^\dagger$ has the identical reduced state on
the system, since the partial trace is invariant under a unitary on the
traced-out factor. Hence $\eta$, built from $\rho$ alone
(Eq.~\eqref{eq:eta-purity}), and any $f(\eta)$, is unchanged. This includes
reading the register (a unitary copying its outcome to a classical
pointer, itself part of $R$) and choosing a measurement basis (a unitary
rotation of $R$ before projection); both are covered by the argument and
checked numerically in Sec.~\ref{sec:checks}. Two of the alternatives
compared are not of this form: sourcing by the measured outcome depends on which
basis the register is measured in (a non-unitary projection, not covered by
the argument above), and sourcing by how many independent copies of a record
exist depends on a count that a unitary redistributing information among
those copies can change.

% =============================================================================
\section{The adapter}
\label{app:adapter}

\subsection{The filtering equations}

For a damped harmonic oscillator under continuous position monitoring at
rate $\mu$ and total momentum diffusion $D$, the steady-state covariance
$(V_x,V_{xp},V_p)$ solves
\begin{align}
  \dot V_x  &= \tfrac{2}{m}V_{xp} - 8\mu V_x^2, \\
  \dot V_{xp} &= \tfrac{1}{m}V_p - m\omega_0^2V_x - \gamma V_{xp} - 8\mu
                 V_xV_{xp}, \\
  \dot V_p  &= -2m\omega_0^2V_{xp} - 2\gamma V_p + 2D - 8\mu V_{xp}^2.
\end{align}
Setting the right-hand sides to zero leaves one quartic in $V_x$ with
positive coefficients, so the positive root is unique. It is found in
$\log V_x$, since the root can lie twenty orders of magnitude below its
$\mu=0$ value. $\mu=D/\hbar^2$ (every channel
conditioned on, at unit efficiency) gives the fully conditioned covariance;
$\mu=\eta_{\det}\Lambda_{\mathrm{ro}}$ (the readout alone) gives the
measurement-conditioned one; $\mu=0$ gives the unconditioned one. The
unconditioned covariance reproduces equipartition, and the fully conditioned
one the ground state.

\subsection{The extrapolation, checked a second way}
\label{app:adapter-check}

The conclusion $f=0$ for a thermal oscillator follows from reading
$\eta_{\mathrm{eff}}$ off the unconditioned state's purity, an extension of
Eq.~\eqref{eq:eta-purity} beyond the two-branch states for which it is
derived.
An independent reading instead asks for the reduced state's coherence at the
material's own lattice localization length $\sigma$,
\begin{equation}
  \eta_{\mathrm{coh}}(\sigma) = \exp\!\left[-\frac{\big(\Sigma-\hbar^2/4V_x
  \big)\,\sigma^2}{2\hbar^2}\right],\qquad \Sigma\equiv V_p-\frac{V_{xp}^2}
  {V_x},
\end{equation}
which shares no algebra with the purity-based route. On both hot platforms
(osmium, torsion balance) the thermal de Broglie length is $10^{-16}\,
\mathrm{m}$ or smaller against a lattice localization of a few picometres, so
$\eta_{\mathrm{coh}}(\sigma)\to0$, and the conclusion $f=0$ does not depend
on which reading is used. % Claim 20

% =============================================================================
\section{Platform parameters}
\label{app:platforms}

\begin{table*}[t]
  \centering
  \caption{All three platforms' parameters, with sources and assumptions
    (expanding Table~\ref{tab:platforms}).}
  \label{tab:platformsfull}
  \input{../../build/latex/tables/platforms-full.tex}
\end{table*}

% =============================================================================
\section{The closed-form resolvability, and its regime of validity}
\label{app:resolvability}

Equation~\eqref{eq:resolvability} assumes two things: weak measurement
($4m\Lambda\hbar\ll(m\omega_0)^2$, so the fully conditioned variance is the
ground state's), and a small shift ($\omega_{\mathrm{SN}}\ll\omega_0$, so the
offset is $\omega_{\mathrm{SN}}^2/2\omega_0$). Checked against the full
filtering model, it is accurate to within a few per cent on the osmium
platform and within tens of per cent on the ground-state-cooled sphere, but
wrong by six orders of magnitude on the torsion balance, which is deep in
the strong-measurement regime and has $\omega_{\mathrm{SN}}$ above
$\omega_0$. All quoted numbers come from the full filtering model;
Eq.~\eqref{eq:resolvability} is used only for scaling arguments, where it
holds.

\bibliography{../../build/latex/refs}
\bibliographystyle{apsrev4-2}

\end{document}
